\begin{enunciado}{\ejExtra[23/06/26 🌶Ejercicio de la Práctica]}
  \fechaEjercicio{23/06/26 Ejercicio de la Práctica}

  Calcular la cantidad de funciones inyectivas $f: \set{1,2,\ldots, 7} \to \set{1,2, \ldots, 11}$
  que cumplen simultáneamente
  \begin{itemize}
    \item $f(1) < f(7)$
    \item $f(3)$ es par
  \end{itemize}
\end{enunciado}

\begin{itemize}
  \item Para cumplir que $f(1) < f(7)$ tomo 2 números cualesquiera del conjunto de llegada,
        la cantidad de formas de elegir esos dos valores sería:
        $$
          \binom{11}{2}.
        $$
        Eso me resuelve la primera condición. Recordar que contar con el número combinatorio
        a ciegas \faIcon{blind} funciona lo más bien, porque entre 2 números distintos que yo elija
        siempre van a ser \textit{uno mayor que el otro} y el número combinatorio
        \textit{no cuenta permutaciones}, es decir que, si agarro el \textit{3 y el 2}, los
        ordeno tipo:
        $$
          f(1) = 2 ~~\text{y} ~~ f(7) = 3
        $$
        y listo, \textit{hacerlo al revés no entra dentro del cálculo del número combinatorio}

  \item La segunda condición: \textit{$f(3)$ es par} está condicionada a los números que tomé
        en el paso anterior, dado que no sé si tomé números \ul{pares}, \ul{impares}
        $\ua{\red{\text{o}}}{\text{\red{\atencion}}}$
        \ul{uno de cada}.
        \begin{enumerate}[label=(\faIcon{calculator}$_{\arabic*}$)]
          \item
                \ul{\textit{Si tomé 2 números pares}}:
                \parrafoDestacado{
                  Me quedan para elegir 3 números pares, entonces
                  puedo elegir para

                  \ul{$f(3) \to  \magenta{3}$ posibles números pares}.
                }

                Por lo tanto si hay $5$ números pares en el \textit{codominio}, puedo
                elegir dos de cinco así:
                $$
                  \binom{5}{2}  = 10 ~ \text{maneras posibles}
                $$

          \item
                \ul{\textit{Si tomé 1 número par y 1 impar}}:
                \parrafoDestacado{
                  Me quedan para elegir 4 números pares, entonces
                  puedo elegir para

                  \ul{$f(3) \to \magenta{4}$ posibles números pares}.
                }

                Por lo tanto si hay $5$ números pares y 6 impares en el \textit{codominio}, puedo
                elegir uno de cinco y uno de seis así:
                $$
                  \binom{5}{1} \cdot \binom{6}{1} = 30 ~ \text{maneras posibles}.
                $$

          \item
                \ul{\textit{Si tomé 2 números impares}}:
                \parrafoDestacado{
                  Me quedan para elegir 5 números pares, entonces
                  puedo elegir para

                  \ul{$f(3) \to \magenta{5}$ posibles números pares}.
                }
                Por lo tanto si hay 6 impares en el \textit{codominio}, puedo
                elegir dos de seis así:
                $$
                  \binom{6}{2} = 15 ~ \text{maneras posibles}.
                $$
        \end{enumerate}

        Juntando esa info y luego completando de forma \textit{inyectiva} tengo:
        $$
          (10 \cdot 3 + 30 \cdot 4 + 15 \cdot 5) \cdot \frac{8}{(8-4)!} =
          \cajaResultado{
            378.000
          }
        $$
\end{itemize}

\begin{aportes}[2]
  \item \aporte{\dirRepo}{naD GarRaz \github}
  \item \aporte{https://github.com/GonsaAvalos}{Gon. Avalos \github}
\end{aportes}
